\documentclass[10pt,a4paper]{amsart}
\usepackage[margin=19mm]{geometry}
\usepackage{amsmath,amssymb,mathtools}
\usepackage[scr=rsfs,cal=euler]{mathalfa}
\usepackage{xfrac}
\usepackage[unicode,hidelinks,bookmarks=false]{hyperref}
\newcommand{\ie}{i.e.\ }
\newcommand{\cf}{cf.\ }
\newcommand{\resp}{respectively\ }
\newcommand{\Subsection}{\subsection}
\providecommand{\nobreakdash}{\nobreak\hbox{-}}
\setlength{\emergencystretch}{2em}
\multlinegap0pt
\usepackage{enumerate}
\usepackage{xcolor}
\usepackage{amssymb}
\usepackage[shortlabels]{enumitem}
\newtheorem{lemma}{Lemma}
\newtheorem{proposition}{Proposition}
\newtheorem{remark}{Remark}
\newtheorem{theorem}{Theorem}
\newcommand{\sub}{\subseteq}
\newcommand{\zN}{\mathbb N}
\title{Frequently recurrent backward shifts}
\author{Rodrigo Cardeccia, Santiago Muro}
\date{Source: arXiv:2407.11799v2 (CC BY 4.0)}
\begin{document}
\maketitle

\noindent\emph{Source and licence.} Frequently recurrent backward shifts. Version arXiv:2407.11799v2 (CC BY 4.0), distributed by arXiv under the Creative Commons Attribution 4.0 International licence, http://creativecommons.org/licenses/by/4.0/, as recorded in the arXiv licence metadata for this exact version and shown on the abstract page https://arxiv.org/abs/2407.11799v2. Full licence notice: \texttt{licenses/arxiv-2407.11799-CC-BY-4.0.txt}.\par\smallskip

\noindent\textit{Editorial scope.} Counted unit: the article's theorem recurrence (source0.tex lines 660--662). The article's complete original proof of it is delivered with it (source0.tex lines 664--759, one \texttt{proof} environment of 96 lines). No mathematics is written, repaired or re-derived: the statement and the proof are the article's own lines, in the article's own notation, and no retained line is substituted or re-flowed.  Retained uncounted as the source-local context the proof needs: lines 219--219; lines 225--227; lines 258--260; lines 593--601; lines 615--617.  Nothing was omitted from any retained span, so the package carries no omission record.  No retained line contains a citation, so the wrapper carries no bibliography and no bibliography entry.  No retained line contains a figure environment, an image-inclusion command or drawing code, so no image is delivered and the package declares no asset dependency.  Editorial formatting only, in addition: an AMS \texttt{amsart} preamble that restates the article's own declarations and macros used by the retained lines, namely the environments \texttt{lemma}, \texttt{proposition}, \texttt{remark}, \texttt{theorem} on separate counters, together with the standard packages those lines need.  The retained environments print in this document as Lemma 1, Proposition 1, Remark 1, Theorem 1 in order of appearance, whereas the article prints the same items with the numbering of its own full document.  The complete native source of the article as distributed in the arXiv e-print for this exact version is bundled unchanged as \texttt{sources/162411/source0.tex}.
\section{Preliminaries and notation}\label{preliminaries}

\begin{equation}\label{incondicionalidad}
    \|ax\|_p\le C_p\|a\|_{\infty} \|x\|_{m_p}.
\end{equation} 

\begin{lemma}[Separation Lemma]\label{separation lemma}
Let $(A_p)_p$ be a sequence of sets with positive lower density and $(k_p)$ a sequence of natural numbers. Then there exist pairwise disjoint sets of positive lower density $(B_p)_p$, such that for every $p$, $B_p\sub A_p$ and such that for every $n\in B_p$, $m\in B_q$, $n\ne m$, we have that $|n-m|>k_p+k_q$.
\end{lemma}

\begin{proposition}\label{prop: condic necesaria vectores recurrentes}
    Let $(X,T)$ be a dynamical system (i.e. $X$ a first-countable topological space and $T:X\to X$  continuous). Let $x\in X$ be a recurrent vector, $(U_p)_p$ be a decreasing base of neighbourhoods of $x$ and $A_p\subset N_T(x,U_p)$ be non-empty sets. Then for each $q\in \zN$ there exists a subsequence $(A_{M_p})_p$ such that 
    % the set defined by the relations
    % \begin{equation}\label{definicion A}
    %     A=\bigcup_{p=0}^\infty\bigcup_{j=0}^p A_{M_p}+N_j ,\, \text{ and, }\,\\
    %  N_k=\min\{ n>N_{k-1}\,:\, n\in\bigcup_{p=0}^{k-1}\bigcup_{j=0}^p A_{M_p}+N_j\}
    %    \end{equation}
       $RAT((A_{M_p})_p)\subseteq N_T(x,U_q)$. 
\end{proposition}

\begin{remark}\label{rem:M_k>N_k}
    Note that in the above proof $M_k$ is chosen after defining $N_k$, so we may assume without loss of generality that $M_k>N_k$.
\end{remark}

\begin{theorem}[Zero-one law for frequently recurrent vectors]\label{0-1 law frequent recurrence}
Let $X$ be a Fr\'echet sequence space with unconditional basis $(e_n)_{n\geq 0}$. If $B$ supports a nontrivial frequently recurrent vector then $B$ is frequently recurrent.
\end{theorem}

\begin{proof}
Let $x\neq 0$ be a frequently recurrent vector. It suffices to prove that for each $R\in\mathbb N_0$, and each $\varepsilon>0$ there is a frequently recurrent vector $y\in X$ such that $\|y-e_R\|<\varepsilon$, where $\|\cdot\|$ denotes an $F$-norm on $X$ (see Section \ref{preliminaries}). 

We may suppose that   $|x_R|>\delta$ for some $\delta>0$.   
Since $x$ is a recurrent vector,  there is for every $p\in\mathbb N_0$, a  $q_p>p+R$ such that $|x_{q_p}|>\delta$. By the continuity of the coordinate functionals,  there are a decreasing sequence of positive numbers $(\varepsilon_p)_p$ and a sequence of natural numbers $(k_p)_p$ such that for each $p\in\mathbb N_0$, $\varepsilon_{p}<\frac{\varepsilon}{2^{p+1}}$, $k_p>p+1$ and   
\begin{align}\label{eq: beta_p (zero one freq) UNIL}
 \begin{split}
  \|w\|_{k_p}<\varepsilon_{p}\Longrightarrow  |w_{j}|<\frac{\delta}{4} ,  \text{ for every } 0\leq j\leq q_p, &\quad \text{and}\\
 \|x-\sum_{0\leq j\le k_p} x_j e_j\|_{m_p}<\frac{\varepsilon}{2^{p+1}}\frac{\delta}{8C_p},&
 \end{split}
 \end{align}
 where $C_p$ and $\|\cdot\|_{m_p}$ are the unconditional constant basis and seminorm associated to $\|\cdot\|_p$ (see \eqref{incondicionalidad}). 



We will now present sets $A_p\sub \zN$  that will allow us to construct a frequently recurrent vector near $e_R$. For each $p\in\mathbb N_0$, $A_p$ will be a set of positive lower  density such that  
\begin{equation}\label{eq: A_p zero one freq}
A_p\sub N_B(x,U_p),\quad\text{where }\, U_p=\left\{y\,:\, \|y-x\|_{k_p+m_p}<\varepsilon_p\frac{\delta}{8C_p}\right\}.
\end{equation}

Applying the Separation Lemma \ref{separation lemma}
we may assume that the $A_p$'s are pairwise disjoint and that for every $n\in A_p,m\in A_q$ with $n\neq m$ we have that
\begin{equation}\label{eq separation zero one freq}
|n-m|>k_p+k_q+2R>p+1+q+1+2R.    
\end{equation}

Notice that for each  $p\in \zN_0$ and each $0\leq j\leq p$ we have that, $\sum_{n\in A_p} e_{n+R+j}$ converges. Indeed, by \eqref{eq: beta_p (zero one freq) UNIL} and \eqref{eq: A_p zero one freq} we have that $|x_{n+q_p}-x_{q_p}|<\frac{\delta}{4}$ and hence $|x_{n+q_p}|>\frac{\delta}{2},
$ for all $n\in A_p.$ By unconditionality this implies that $\sum_{n\in A_p} e_{n+q_p}$ converges and by applying the backward shift $q_p-(R+j)$ times it follows that $\sum_{n\in A_p} e_{n+j+R}$ converges for every $0\leq j\leq p$. By the finite invariance property of  the sets of positive lower density, 
we may assume that 
\begin{equation}\label{eq A_p empieza en p}
    A_p\sub [p+R+1,\infty)
\end{equation}
and that,  for every $0\le j\le p,$
\begin{equation}\label{eq: sumas de e_j norma chica (zero one freq)}
 \|\sum_{n\in A_p} e_{n+R+j}\|<\frac{\varepsilon}{(p+1)2^{p+1}}.   
\end{equation}

Applying Proposition \ref{prop: condic necesaria vectores recurrentes} and Remark \ref{rem:M_k>N_k} there is a sequence $(M_p)_p$ for which $RAT((A_{M_p})_p)\subseteq N_B(x,U_0)$ and such that $M_p>N_p$ for every $p\geq 1$ and $M_0=N_0=0$. The fact that $M_p>N_p$ together with \eqref{eq A_p empieza en p}, implies that $(A_{M_p})_p$ is compatible, i.e. 
 $RAT((A_{M_p})_p)=\{N_k\}$.
 % Let  $(M_p)_p$ be the sequence given by Proposition \ref{prop: condic necesaria vectores recurrentes} (for $q=0$).   Consider the recurrent arithmetic thickening $RAT((A_{M_p})_p)=\bigcup_{p=0}^{\infty}\bigcup _{j=0}^p (A_{M_p}+N_j)$, where $N_{k}=\min_{n>N_{k-1}}\{n\in\bigcup_{p=0}^{k-1}\bigcup _{j=0}^p (A_{M_p}+N_j) \}$. By Remark \ref{rem:M_k>N_k} We may assume that $N_k<M_k$ for every $k\ge1$ which, together with \eqref{eq A_p empieza en p}, implies that $(A_{M_p})_p$ is compatible, i.e. 
 % $RAT((A_{M_p})_p)=\{N_k\}$. 

We will show that the vector
\begin{align*}
y:=\hspace{-.3cm}\sum_{n\in RAT((A_{M_p})_p) }\hspace{-.3cm}e_{n+R}=e_R+\sum_{p=0}^{\infty}\sum_{n\in A_{M_p}}\sum_{j=0}^pe_{n+N_j+R}=e_R+\sum_{p=0}^\infty\sum_{n\in A_{M_p}}\sum_{j=0}^{N_p+R}y_je_{n+j}=\sum_{j=0}^\infty e_{N_j+R}.
\end{align*}
is a (well-defined) frequently recurrent vector such that $\|y-e_R\|<\varepsilon.$ The well-definition and the inequality $\|y-e_R\|<\varepsilon$ follow from \eqref{eq: sumas de e_j norma chica (zero one freq)}. 

We prove that $y$ is frequently recurrent. 
 Fix $q\in\zN$ and $m\in A_{M_{q}}$. We will show that $B^m(y)$ is close to $\sum_{j=0}^{N_{q}+R}y_je_j$.
 
 
 Recall that by \eqref{eq separation zero one freq}, since $k_p>p$ for every $p,$ we have that for every $p\in\mathbb N$ and $n\in A_{M_p}$, $|n-m|>M_p+M_{q}+2R>N_p+N_{q}+2R$. Thus, if $n<m$ and $n\in A_{M_{p}}$, then $n-m+N_j+N_q+R<0$ for every $0\le j\le p$. Hence,
$$B^m\left(\sum_{n\in A_{M_p},n<m}\sum_{j=0}^pe_{n+N_j+R}\right)=\sum_{n\in A_{M_p},n<m}\sum_{j=0}^pe_{n-m+N_j+R}=0.$$
Thus, noting that $m$ belongs only to $A_{M_{q}}$, because the $A_{M_p}$'s are pairwise disjoint, we get that
\begin{align*}
B^m(y)-\sum_{j=0}^{N_{q}+R}y_je_j &=B^m(y)-\sum_{j=0}^{q}e_{N_j+R}\\
&=\sum_{p\neq q}\sum_{n\in A_{M_p},n>m}\sum_{j=0}^{p}e_{n-m+N_j+R}+\sum_{n\in A_{M_{q}},n\geq m}\sum_{j=0}^{q}e_{n-m+N_j+R}-\sum_{j=0}^qe_{N_j+R}\\
&=\sum_{p=0}^\infty \sum_{n\in A_{M_p},n>m}\sum_{j=0}^pe_{n-m+N_j+R}.
\end{align*}
By Proposition \ref{prop: condic necesaria vectores recurrentes},  for $0\le j\le p,$ $n\in A_{M_p}$,
$$
B^{n+N_j}x\in U_0,\quad\text{and in particular,}\quad \| B^{n+N_j}x -x\|_{k_0+m_0}<\varepsilon_0.
$$
Then since $|x_R|>\delta,$ we have by \eqref{eq: beta_p (zero one freq) UNIL} that 
\begin{equation*}
    |x_{n+N_j+R}|\ge |x_R|-|x_{n+N_j+R}-x_R| \ge \delta - \frac{\delta}{4}.
\end{equation*}

Since by \eqref{eq separation zero one freq}, $n-m>k_{M_p}+k_{M_{q}}+R>k_{M_q}>k_{N_q}$ whenever $n\in A_{M_p}$ and $n>m$,  we have  $[\sum_{l=0}^{k_{N_{q}}}x_le_l]_{n-m+N_j+R}=0$ for each $0\le j\le p$. Hence it follows that  
 % \begin{align}\label{cota B^my, n>m}
 % \|B^m(y)-\sum_{j=1}^{N_{s_0}+R}y_je_j \|&\leq CR\|B\|^R\|z\|_\infty\left\|\sum_{s=0}^\infty \sum_{n\in A_{M_s},n>m}\sum_{l=0}^se_{n-m+N_l+R}\right\|\\
 % &\leq \frac{4C^2}{\delta}R\|B\|^R\|z\|_\infty\left\|\sum_{s=1}^{\infty}\sum_{n\in A_{M_s},n>m}\sum_{l=0}^sx_{n+N_l+R}e_{n-m+N_l+R}\right\|\\
 % &=\frac{4C^2}{\delta}R\|B\|^R\|z\|_\infty\left\|\sum_{s=0}^{\infty}\sum_{n\in A_{M_s},n>m}\sum_{l=0}^s[B^m(x)-\sum_{j=1}^{k_{N_{s_0}}}x_j]_{n-m+N_l+R}e_{n-m+N_l+R}\right\|\\
 % &\leq \frac{4C^3}{\delta}R\|B\|^R\|z\|_\infty\left(\|B^m(x)-x\|+\|x-\sum_{j=1}^{k_{N_{s_0}}}x_j\|\right)< \varepsilon_{N_{s_0}}.
 % \end{align}
  \begin{align}\label{eq: B^my final}
 \begin{split}
 \left\| B^m(y)-\sum_{j=0}^{N_{q}+R}y_je_j\right\|_q&=
 \left\|\sum_{p=0}^\infty \sum_{n\in A_{M_p},n>m}\sum_{j=0}^pe_{n-m+N_j+R}\right\|_q\\
 &=\left\|\sum_{p=0}^{\infty}\sum_{n\in A_{M_p},n>m}\sum_{j=0}^p\frac{1}{x_{n+N_j+R}}\cdot{x_{n+N_j+R}}e_{n-m+N_j+R}\right\|_q\\
 &=\left\|\sum_{p=0}^{\infty}\sum_{n\in A_{M_p},n>m}\sum_{j=0}^p\frac{1}{x_{n+N_j+R}}\cdot[B^m(x)-\sum_{l=0}^{k_{N_{q}}}x_le_l]_{n-m+N_j+R}e_{n-m+N_j+R}\right\|_q\\
 &\leq \frac{4C_q}{\delta}\left(\|B^m(x)-x\|_{m_q}+\|x-\sum_{l=0}^{k_{N_{q}}}x_le_l\|_{m_q}\right)\\
 & 
\le \frac{4C_q}{\delta}\left(\|B^m(x)-x\|_{k_{M_q}+m_{M_q}}+\|x-\sum_{l=0}^{k_{N_{q}}}x_le_l\|_{m_{N_q}}\right)< \frac{\varepsilon}{2^{N_{q}+1}},
  \end{split}
 \end{align}
where in the last inequality we have used \eqref{eq: beta_p (zero one freq) UNIL} and the definition of $A_{M_p}$ in \eqref{eq: A_p zero one freq}.


 
Finally let $V$ be an open set around $y$. Then since  ${N_q}\to\infty$ as $q\to\infty,$    $\sum_{j=0}^{N_q+R} y_je_j\to y$ and hence,  for big enough  $q\in\mathbb N$, it follows by \eqref{eq: B^my final} that $A_{M_{q}}\sub N_B(y,V).$
This shows that $y$ is a frequently recurrent vector.


\end{proof}

\end{document}
